\documentclass[11pt]{article}
\usepackage{amsmath,amssymb,amsthm,mathtools}
\usepackage[margin=1in]{geometry}
\theoremstyle{plain}
\newtheorem{theorem}{Theorem}
\newtheorem{lemma}{Lemma}
\theoremstyle{remark}
\newtheorem{remark}{Remark}
\newcommand{\E}{\mathbb{E}}\newcommand{\PP}{\mathbb{P}}\newcommand{\R}{\mathbb{R}}\newcommand{\Z}{\mathbb{Z}}

\title{Formalization brief: de-opaquing the two-particle ASEP Green's function\\
(Sch\"utz's reflection formula in the $F$-basis, and the $C/(1+t)$ decay)}
\date{6 July 2026}

\begin{document}
\maketitle

\noindent\textbf{Goal.} Discharge \texttt{asepGreen\_integral\_decay} in
\texttt{TypeDDecouplingLCLT.lean} (attached). As encoded, the integrand is
\texttt{opaque} --- the \texttt{sorry} is unprovable by design. The task is
therefore a sanctioned RESTRUCTURING of the \texttt{lem\_asep} block (the
opaque \texttt{asepGreenIntegrand}, the \texttt{hGreen} hypothesis,
\texttt{asepGreen\_integral\_decay}, and \texttt{lem\_asep}), preserving or
strengthening \texttt{lem\_asep}'s conclusion
($\mathrm{asepKernel}\ q\ t\ \xi\ \xi'\le C/(1+t)$) and keeping every consumer
(\texttt{thm\_kernel} etc.) elaborating. All repairs documented in docstrings
(fidelity-repair convention). New library-clean file
\texttt{TypeDDecouplingBethe.lean} encouraged. Whole project must build;
standard axioms only.

\section*{Design: work in the $F$-basis, not the contour integral}

Do NOT define a contour integrand. Sch\"utz's two-particle solution for the
partially asymmetric exclusion process has an equivalent representation as a
geometric combination of single-particle kernels (``dynamical reflection''),
which is the formalization-friendly form: no complex analysis, no $S$-matrix
poles, no orthogonality-of-characters subtleties.

\begin{itemize}
\item[(1)] \underline{Single-particle kernel.} Let $F_m(t)$ be the free
asymmetric walk kernel (right rate $r_R$, left rate $r_L$ --- match the
project's same-species dual rates; carry them symbolically). Define it
concretely EITHER by the lattice Fourier integral
$F_m(t)=\frac1{2\pi}\int_{-\pi}^{\pi}
e^{t(r_Re^{\mathrm i\theta}+r_Le^{-\mathrm i\theta}-r_R-r_L)}
e^{-\mathrm im\theta}\,d\theta$ (real part; reuse the inversion machinery of
\texttt{TypeDDecouplingKR.lean}, attached) OR as
$\exp(tA_1)\delta_0$ via \texttt{TypeDDecouplingSemigroup.lean} (attached),
proving the two agree by the uniqueness argument. Prove the basic $F$-calculus:
the forward ODE $\dot F_m=r_RF_{m-1}+r_LF_{m+1}-(r_R+r_L)F_m$, the initial
condition $F_m(0)=\delta_{m,0}$, nonnegativity, $\sum_mF_m=1$, and the decay
$F_m(t)\le C_1/\sqrt{1+t}$ with $C_1=C_1(r_R+r_L)$ (modulus bound
$+\ 1-\cos\theta\ge\tfrac2{\pi^2}\theta^2$, as in the \texttt{lem\_Rlclt}
campaign's Tier 2).
\item[(2)] \underline{The two-particle ansatz.} For the ordered pair
$x_1<x_2$ started from $y_1<y_2$, make the reflection ansatz
\[
  G_t(x_1,x_2)\ =\ F_{x_1-y_1}(t)F_{x_2-y_2}(t)
  \ +\ \sum_{k\ge0}c_k\,F_{x_1-y_2-a_k}(t)\,F_{x_2-y_1+b_k}(t),
\]
with coefficients $c_k$ (geometric in $r_L/r_R$ or $r_R/r_L$) and integer
shifts $a_k,b_k$ TO BE DERIVED: impose the exclusion boundary condition ---
the two-particle master equation for the exclusion pair differs from the free
one only at contact $x_2=x_1+1$, and the standard reduction replaces it by the
free equation plus the boundary identity
\[
  r_R\,G_t(x,x)+r_L\,G_t(x+1,x+1)\ =\ (r_R+r_L)\,G_t(x,x+1)
  \qquad(\forall x,t),
\]
(check the exact form against the project's two-particle rates; derive rather
than trust this template). Substituting the ansatz turns the boundary identity
into a two-term recursion for $(c_k,a_k,b_k)$, solvable in closed form with
$|c_k|\le\rho^k$ for a ratio $\rho=\rho(q)<1$ (this is where $q\in(0,1)$
enters). Verify the derived formula satisfies: (i) the free equation off
contact (from the $F$-ODE, term by term; justify the interchange with the sum
by uniform geometric domination); (ii) the boundary identity (the recursion,
by construction); (iii) the initial condition $G_0=\delta$ on the ordered
sector (from $F_m(0)=\delta_{m,0}$: check the reflection terms vanish at
$t=0$ for $x_1<x_2$ --- the shifts $a_k,b_k$ make the indices nonzero there;
if a low-$k$ term needs a separate case, handle it).
\item[(3)] \underline{Identification.} Define (or, if \texttt{asepKernel} is
opaque, REPLACE it by --- fidelity repair, documented) the two-particle
exclusion kernel via \texttt{TypeDDecouplingSemigroup.lean}: the concrete
finite-range bounded generator on the ordered sector
$\{(x_1,x_2):x_1<x_2\}$. Prove $G=\,$kernel by the weighted-Gronwall
uniqueness (both solve the same forward ODE with the same initial condition;
the boundary reduction of (2) is exactly what makes $G$ a solution of the
EXCLUSION equation).
\item[(4)] \underline{Decay.} From (1)'s $F$-bound and $\sum_k|c_k|\le
(1-\rho)^{-1}$:
$G_t(x_1,x_2)\le C_1^2\big(1+(1-\rho)^{-1}\big)/(1+t)$ --- the target
$C(q)/(1+t)$, with $C$ quantified before $t,\xi,\xi'$ as in the current
statement.
\item[(5)] \underline{Interface.} Restate the block: \texttt{lem\_asep}
becomes UNCONDITIONAL (no \texttt{hGreen} hypothesis --- the formula is now a
theorem); \texttt{asepGreen\_integral\_decay} either becomes the proved decay
of the concrete representation or is retired in favour of it (keep a
deprecation docstring pointing to the new theorem so the paper's \S1.5
mapping stays traceable); the opaque \texttt{asepGreenIntegrand} may be
deleted or kept with a docstring noting it was superseded. Thread consumers.
\end{itemize}

\begin{remark}[fallbacks and hygiene]
(1) Sanctioned fallbacks, in order: if the $t=0$ vanishing of reflection terms
resists in full generality, a case split or an explicit finite verification of
the lowest terms is fine; if the identification (3) resists at the
ordered-sector boundary, the single named hypothesis may carry ONLY the
sector-boundary regularity fact (documented) --- (1),(2),(4) must be proved.
(2) The exact rates: read them off the project's encoding of the same-species
dual pair (do not assume $r_R+r_L=1$); if the encoding leaves them abstract,
carry $(r_R,r_L)$ with $0<r_L<r_R$ (or as encoded) and expose $q$ through
their ratio as the project does.
(3) Sanity identities worth proving as tests: total mass $\le1$;
$G\ge0$; symmetry under the appropriate reflection.
(4) Report: the derived $(c_k,a_k,b_k)$ in the summary; whether
\texttt{lem\_asep} is now unconditional; final term-level count
(target 10 $\to$ 9).
\end{remark}

\end{document}
